% GENERATED FILE. Edit canonical lessons, then run npm run generate:books.
% canonical-source-hash: fnv1a64:a9146fb696381ddd
% canonical-lessons: BN-C390-sorok, BN-C390-bak, BN-C390-gorto, BN-C390-pich, BN-C390-gorano

\chapter{Highways and Bends}
\label{ch:bn-highways-and-bends}
\hlchaptermodality{390}

\begin{chapteropening}
\textbf{By the end of this chapter:} \emph{I can describe a highway, its bends, its tar and its potholes.}

Complete the last lesson of chapter 390: I can describe a highway, its bends, its tar and its potholes.
\end{chapteropening}

\section[\texorpdfstring{\bn{সড়ক} (sôṛôk)}{সড়ক (sôṛôk)}]{\bn{সড়ক} (sôṛôk) — a highway; a main road}

\label{lesson:BN-C390-sorok}



\begin{quote}
Before the new one: say the Bengali for to signal, then the Bengali for to call out.
\end{quote}

\subsection*{You'll want to know: \bn{সড়ক}}
\textbf{\bn{সড়ক}} — \emph{sôṛôk} — \textquotedblleft{}a highway; a main road\textquotedblright{}.

\begin{quote}
\bn{নতুন} \bn{সড়কে} \bn{বাস} \bn{দ্রুত} \bn{চলে।}
\end{quote}

\begin{quote}
\emph{nôtun sôṛôke bās drutô chôle.}
\end{quote}

\begin{quote}
Buses go fast on the new highway.
\end{quote}

\bn{রাস্তা} is any road; \bn{সড়ক} is a main road between towns, the word on maps: \bn{জাতীয়} \bn{সড়ক}, a national highway.

\begin{cousinweb}
\textbf{\bn{সড়ক}} (\emph{sôṛôk}) probably belongs with Sanskrit \emph{sṛtí}, \textquotedblleft{}a road, a path\textquotedblright{}, the entry where CDIAL sets Bengali \emph{saṛak}, \textquotedblleft{}a road\textquotedblright{} (CDIAL 13577); the Hindi word is filed under a second entry as well (CDIAL 12269), so the link is hedged.
\end{cousinweb}

\begin{grammarlens}[title={What you've built}]
One more word for getting around.
\end{grammarlens}

\subsection*{Your turn}
\begin{itemize}
\raggedright
  \item \emph{Say it:} \emph{sôṛôk}
  \item \emph{Say it:} \emph{sôṛôk}, once more
  \item \emph{Say it:} \emph{nôtun sôṛôke bās drutô chôle.} — Buses go fast on the new highway.
  \item \emph{Recall:} say the Bengali for to signal, then the Bengali for to call out, then say \emph{sôṛôk} again
\end{itemize}

\begin{tcolorbox}[breakable,colback=teal!4,colframe=teal!35!black,title={Before you move on}]
What does \bn{সড়ক} mean? (\textquotedblleft{}A highway; a main road\textquotedblright{}.) Say it once more.
\end{tcolorbox}
\section[\texorpdfstring{\bn{বাঁক} (bā̃k)}{বাঁক (bā̃k)}]{\bn{বাঁক} (bā̃k) — a bend; a turn (in a road or river)}

\label{lesson:BN-C390-bak}



\begin{quote}
Before the new one: say the Bengali for to call out, then the Bengali for a highway.
\end{quote}

\subsection*{You'll want to know: \bn{বাঁক}}
\textbf{\bn{বাঁক}} — \emph{bā̃k} — \textquotedblleft{}a bend; a turn (in a road or river)\textquotedblright{}.

\begin{quote}
\bn{সামনের} \bn{বাঁকে} \bn{বাসটা} \bn{ধীরে} \bn{চলল।}
\end{quote}

\begin{quote}
\emph{sāmner bā̃ke bāsṭā dhīre chôllô.}
\end{quote}

\begin{quote}
At the bend ahead the bus slowed down.
\end{quote}

\bn{বাঁকা}, crooked, is the same root as an adjective.

\begin{cousinweb}
\textbf{\bn{বাঁক}} (\emph{bā̃k}) is inherited: CDIAL gives \emph{vaṅka}, \textquotedblleft{}bent, crooked\textquotedblright{}, with the sense \textquotedblleft{}a bend of a river\textquotedblright{}, and lists Bengali \emph{bā̃kā}, \textquotedblleft{}crooked\textquotedblright{}, under it (CDIAL 11191); the nasal of \emph{vaṅ} survives as the \bn{ঁ}.
\end{cousinweb}

\begin{grammarlens}[title={What you've built}]
One more word for getting around.
\end{grammarlens}

\subsection*{Your turn}
\begin{itemize}
\raggedright
  \item \emph{Say it:} \emph{bā̃k}
  \item \emph{Say it:} \emph{bā̃k}, once more
  \item \emph{Say it:} \emph{sāmner bā̃ke bāsṭā dhīre chôllô.} — At the bend ahead the bus slowed down.
  \item \emph{Recall:} say the Bengali for to call out, then the Bengali for a highway, then say \emph{bā̃k} again
\end{itemize}

\begin{tcolorbox}[breakable,colback=teal!4,colframe=teal!35!black,title={Before you move on}]
What does \bn{বাঁক} mean? (\textquotedblleft{}A bend; a turn (in a road or river)\textquotedblright{}.) Say it once more.
\end{tcolorbox}
\section[\texorpdfstring{\bn{গর্ত} (gôrtô)}{গর্ত (gôrtô)}]{\bn{গর্ত} (gôrtô) — a hole; a pothole}

\label{lesson:BN-C390-gorto}



\begin{quote}
Before the new one: say the Bengali for a highway, then the Bengali for a bend.
\end{quote}

\subsection*{You'll want to know: \bn{গর্ত}}
\textbf{\bn{গর্ত}} — \emph{gôrtô} — \textquotedblleft{}a hole; a pothole\textquotedblright{}.

\begin{quote}
\bn{গাড়িটা} \bn{একটা} \bn{গর্তে} \bn{পড়ল।}
\end{quote}

\begin{quote}
\emph{gāṛiṭā ekṭā gôrte pôṛlô.}
\end{quote}

\begin{quote}
The car went into a pothole.
\end{quote}

\begin{cousinweb}
\textbf{\bn{গর্ত}} (\emph{gôrtô}) is Sanskrit \emph{garta}, \textquotedblleft{}a hollow, a hole, a pit\textquotedblright{} (Monier-Williams), taken over whole as a tatsama.
\end{cousinweb}

\begin{grammarlens}[title={What you've built}]
One more word for getting around.
\end{grammarlens}

\subsection*{Your turn}
\begin{itemize}
\raggedright
  \item \emph{Say it:} \emph{gôrtô}
  \item \emph{Say it:} \emph{gôrtô}, once more
  \item \emph{Say it:} \emph{gāṛiṭā ekṭā gôrte pôṛlô.} — The car went into a pothole.
  \item \emph{Recall:} say the Bengali for a highway, then the Bengali for a bend, then say \emph{gôrtô} again
\end{itemize}

\begin{tcolorbox}[breakable,colback=teal!4,colframe=teal!35!black,title={Before you move on}]
What does \bn{গর্ত} mean? (\textquotedblleft{}A hole; a pothole\textquotedblright{}.) Say it once more.
\end{tcolorbox}
\section[\texorpdfstring{\bn{পিচ} (pich)}{পিচ (pich)}]{\bn{পিচ} (pich) — tar (on a road)}

\label{lesson:BN-C390-pich}



\begin{quote}
Before the new one: say the Bengali for a bend, then the Bengali for a pothole.
\end{quote}

\subsection*{You'll want to know: \bn{পিচ}}
\textbf{\bn{পিচ}} — \emph{pich} — \textquotedblleft{}tar (on a road)\textquotedblright{}.

\begin{quote}
\bn{গরমে} \bn{রাস্তার} \bn{পিচ} \bn{নরম} \bn{হয়ে} \bn{গেল।}
\end{quote}

\begin{quote}
\emph{gôrôme rāstār pich nôrôm hôye gelô.}
\end{quote}

\begin{quote}
In the heat the tar on the road went soft.
\end{quote}

\bn{পিচের} \bn{রাস্তা}, a tarred road; a \bn{কাঁচা} \bn{রাস্তা} is a dirt road.

\begin{cousinweb}
\textbf{\bn{পিচ}} (\emph{pich}) is borrowed from English \emph{pitch}, the black tar that roads are sealed with.
\end{cousinweb}

\begin{grammarlens}[title={What you've built}]
One more word for getting around.
\end{grammarlens}

\subsection*{Your turn}
\begin{itemize}
\raggedright
  \item \emph{Say it:} \emph{pich}
  \item \emph{Say it:} \emph{pich}, once more
  \item \emph{Say it:} \emph{gôrôme rāstār pich nôrôm hôye gelô.} — In the heat the tar on the road went soft.
  \item \emph{Recall:} say the Bengali for a bend, then the Bengali for a pothole, then say \emph{pich} again
\end{itemize}

\begin{tcolorbox}[breakable,colback=teal!4,colframe=teal!35!black,title={Before you move on}]
What does \bn{পিচ} mean? (\textquotedblleft{}Tar (on a road)\textquotedblright{}.) Say it once more.
\end{tcolorbox}
\section[\texorpdfstring{\bn{গড়ানো} (gôṛāno)}{গড়ানো (gôṛāno)}]{\bn{গড়ানো} (gôṛāno) — to roll; to roll down}

\label{lesson:BN-C390-gorano}



\begin{quote}
Before the new one: say the Bengali for a pothole, then the Bengali for tar.
\end{quote}

\subsection*{You'll want to know: \bn{গড়ানো}}
\textbf{\bn{গড়ানো}} — \emph{gôṛāno} — \textquotedblleft{}to roll; to roll down\textquotedblright{}.

\begin{quote}
\bn{পাথরটা} \bn{পাহাড়} \bn{থেকে} \bn{গড়িয়ে} \bn{পড়ল।}
\end{quote}

\begin{quote}
\emph{pāthôrṭā pāhāṛ theke gôṛiye pôṛlô.}
\end{quote}

\begin{quote}
The stone rolled down the hill.
\end{quote}

Past (he/she …): \bn{সেই} \bn{পাথর} \bn{গড়াল}, \emph{sei pāthôr gôṛālô}.

Not \bn{গড়া}, to build: a different verb that sounds close.

\begin{cousinweb}
\textbf{\bn{গড়ানো}} (\emph{gôṛāno}) is inherited: CDIAL lists Bengali \emph{gaṛāna}, \textquotedblleft{}to make slide or roll\textquotedblright{}, under Sanskrit \emph{gáḍati}, \textquotedblleft{}drops, runs as a liquid\textquotedblright{} (CDIAL 3973).
\end{cousinweb}

\begin{grammarlens}[title={What you've built}]
One more word for getting around.
\end{grammarlens}

\subsection*{Your turn}
\begin{itemize}
\raggedright
  \item \emph{Say it:} \emph{gôṛāno}
  \item \emph{Say it:} \emph{gôṛāno}, once more
  \item \emph{Say it:} \emph{pāthôrṭā pāhāṛ theke gôṛiye pôṛlô.} — The stone rolled down the hill.
  \item \emph{Recall:} say the Bengali for a pothole, then the Bengali for tar, then say \emph{gôṛāno} again
\end{itemize}

\begin{tcolorbox}[breakable,colback=teal!4,colframe=teal!35!black,title={Before you move on}]
What does \bn{গড়ানো} mean? (\textquotedblleft{}To roll; to roll down\textquotedblright{}.) Say it once more.
\end{tcolorbox}
